\sectionkicker{Source-backed technical notes}
\subsection*{Paper Watch — 長時間Agentを測る・守る: Technical Notes}
本欄は記事本文で圧縮した一次資料上の情報を、比較・再検証しやすい形で整理したものである。時系列、確認済みの主張、留意点、一次資料URLを掲載する。
\medskip
\subsection*{Theme at a glance}
\addcontentsline{toc}{subsection}{Theme at a glance}
\begin{center}
\footnotesize
\begin{tabularx}{\linewidth}{@{}p{0.20\linewidth}p{0.10\linewidth}p{0.14\linewidth}X@{}}
\toprule
資料 & 位置づけ & 種別 & 時系列 \\
\midrule
The Unfireable Safety Kernel: Execution-Time AI Alignment for AI Agents and Other Escapable AI Systems & 主要資料 & 論文 & 2026-06-24 (論文\_RELEASE) \\
Diagnosing and Mitigating Context Rot in Long-horizon Search & 主要資料 & 論文 & 2026-06-29 (論文\_RELEASE) \\
MemDelta: Controlled Baselines and Hidden Confounds in Agent Memory Evaluation & 主要資料 & 論文 & 2026-06-29 (論文\_RELEASE) \\
\bottomrule
\end{tabularx}
\end{center}
\normalsize

\begin{technicalnote}{The Unfireable Safety Kernel: Execution-Time AI Alignment for AI Agents and Other Escapable AI Systems}{主要資料}
\begin{tabularx}{\linewidth}{@{}>{\bfseries}p{0.22\linewidth}X@{}}
組織 & Paper authors \\
種別 & 論文 \\
時系列 & 2026-06-24 (論文\_RELEASE) \\
\end{tabularx}
\smallskip
{\bfseries 一次資料から整理したtechnical points}
\begin{itemize}[leftmargin=1.5em,itemsep=0.35em]
\item \textbf{Author claim}: The authors propose an execution-time safety kernel that remains outside an agent's own mutable runtime.
\end{itemize}
{\bfseries 読む際の境界}
\begin{itemize}[leftmargin=1.5em,itemsep=0.35em]
\item \textbf{分析上の留意点}: The evidence is the authors' arXiv paper/model card and reported evaluations; results are not independently reproduced in this Evidence review.
\end{itemize}
{\bfseries 一次資料}
\begingroup\sloppy
\begin{itemize}[leftmargin=1.5em,itemsep=0.25em]
\item {\scriptsize\url{http://arxiv.org/abs/2606.26057v1}}
\end{itemize}
\endgroup
\end{technicalnote}
\medskip

\begin{technicalnote}{Diagnosing and Mitigating Context Rot in Long-horizon Search}{主要資料}
\begin{tabularx}{\linewidth}{@{}>{\bfseries}p{0.22\linewidth}X@{}}
組織 & Paper authors \\
種別 & 論文 \\
時系列 & 2026-06-29 (論文\_RELEASE) \\
\end{tabularx}
\smallskip
{\bfseries 一次資料から整理したtechnical points}
\begin{itemize}[leftmargin=1.5em,itemsep=0.35em]
\item \textbf{Author claim}: The authors identify premature termination as a long-context failure mode in long-horizon search.
\end{itemize}
{\bfseries 読む際の境界}
\begin{itemize}[leftmargin=1.5em,itemsep=0.35em]
\item \textbf{分析上の留意点}: The evidence is the authors' arXiv paper/model card and reported evaluations; results are not independently reproduced in this Evidence review.
\end{itemize}
{\bfseries 一次資料}
\begingroup\sloppy
\begin{itemize}[leftmargin=1.5em,itemsep=0.25em]
\item {\scriptsize\url{http://arxiv.org/abs/2606.29718v2}}
\end{itemize}
\endgroup
\end{technicalnote}
\medskip

\begin{technicalnote}{MemDelta: Controlled Baselines and Hidden Confounds in Agent Memory Evaluation}{主要資料}
\begin{tabularx}{\linewidth}{@{}>{\bfseries}p{0.22\linewidth}X@{}}
組織 & Paper authors \\
種別 & 論文 \\
時系列 & 2026-06-29 (論文\_RELEASE) \\
\end{tabularx}
\smallskip
{\bfseries 一次資料から整理したtechnical points}
\begin{itemize}[leftmargin=1.5em,itemsep=0.35em]
\item \textbf{Author claim}: The authors show that uncontrolled model, embedding, and retrieval changes can confound comparisons between agent-memory systems.
\end{itemize}
{\bfseries 読む際の境界}
\begin{itemize}[leftmargin=1.5em,itemsep=0.35em]
\item \textbf{分析上の留意点}: The evidence is the authors' arXiv paper/model card and reported evaluations; results are not independently reproduced in this Evidence review.
\end{itemize}
{\bfseries 一次資料}
\begingroup\sloppy
\begin{itemize}[leftmargin=1.5em,itemsep=0.25em]
\item {\scriptsize\url{http://arxiv.org/abs/2606.29914v1}}
\end{itemize}
\endgroup
\end{technicalnote}
\medskip
